\documentclass{situationist-preprint}

\title{Practical Rationality and Lottery Beliefs}
\msid{EPI-09-b}
\msfamily{EPI-09}
\msdiscipline{Epistemology}
\mstype{Companion manuscript}
\mskeywords{lottery; practical reasoning; knowledge norm; pragmatic encroachment; acceptance}

\begin{document}
\maketitle

\begin{abstract}
EPI-09-a argues that the holder of a lottery ticket does not know that it will lose but may rationally act as if it will. We develop the practical half of the thesis. The criterion of EPI-09-a explains the case that troubles it most, the offer to buy the ticket for a small sum, and shows that the threshold at which acting as if becomes irrational is the ticket's expected value. We argue that the variation of this threshold with stakes is no ground for pragmatic encroachment on knowledge, since what varies is the permissibility of an act, not the subject's epistemic position; and that acting as if a proposition is true requires accepting it for the purposes of a decision, not believing it outright.
\end{abstract}

\section{Selling the ticket}

The holder plans her affairs on the assumption that her ticket will lose: she does not book the safari. Yet if someone offers her a penny for the ticket, she should not sell, although selling is what she would do if she were sure it would lose \citep{Hawthorne2004}. The criterion of EPI-09-a separates the two cases. An agent acts as if $p$ when she chooses an act that would be best were $p$ true (its Definition~2.3), and by Proposition~3.2 of EPI-09-a acting as if her ticket loses, in a decision between an act $a$ best if it loses and an alternative $b$, is permitted if and only if $k \leq (N-1)g$, where $g$ is the advantage of $a$ if the ticket loses and $k$ the advantage of $b$ if it wins.

\begin{proposition}\label{p:sell}
Let the prize be $P$ and the offered price $s$, with $0 < s < P$, in a lottery of $N$ tickets. Acting as if the ticket will lose, by selling it, is rationally permitted if and only if $s \geq P/N$.
\end{proposition}

\begin{proof}
Selling is best if the ticket loses, with advantage $g = s$; keeping is best if it wins, with advantage $k = P - s$. By Proposition~3.2 of EPI-09-a, selling is permitted if and only if $P - s \leq (N - 1)s$, that is, $P \leq Ns$.
\end{proof}

The holder may act as if her ticket will lose when the act is declining an expensive trip, and may not when the act is selling below expected value. In both decisions her credence and her evidence are the same; what differs is the ratio of what she stands to lose to what she stands to gain.

\section{Stakes and knowledge}

\citet{FantlMcGrath2002} argue that whether one's evidence is good enough to justify believing $p$, and so whether one can know $p$, depends on whether one is rational to act as if $p$; since the latter varies with the stakes, so does the former. The lottery shows that the variation in what one may rationally do can be explained without such encroachment. By Theorem~3.1 of EPI-09-a the holder does not know in either decision, because a world in which her ticket wins is close; her epistemic position is the same in both. Yet she is rational to act as if her ticket will lose in the first decision and not in the second, and Proposition~\ref{p:sell} explains the difference by expected utility alone. Rationally acting as if $p$ therefore does not require knowing $p$, and the difference between the decisions is not a difference in knowledge. Stakes vary the permissibility of acts on a fixed epistemic position; they need not vary the position.

\begin{objection}
Practical reasoning that uses \enquote{my ticket will lose} as a premise treats it as a reason, and one may treat as reasons only what one knows \citep{HawthorneStanley2008}.
\end{objection}

\begin{reply}
The holder's practical reasoning need not use that premise. Its premises are the odds and the payoffs, which she knows, and its conclusion, that the act best if the ticket loses has the greater expected utility, follows from them. The proposition that the ticket will lose figures in her deliberation as a working assumption that simplifies planning, adopted because adopting it is licensed by what she knows. In the sense of \citet{Bratman1992}, she accepts it in the context of the decision, and she drops the assumption, without any change in belief, when a decision such as the offer for the ticket makes the ratio unfavorable.
\end{reply}

\section{Conclusion}

The holder does not know that her ticket will lose, and need not believe it outright. She may treat its losing as settled for the purposes of a decision exactly when the expected-utility condition of EPI-09-a holds, which in the case of a sale is the condition that the price be at least the ticket's expected value. The lottery thus separates the norms of action from the norm of knowledge without making either depend on the stakes of the other.

\begin{thebibliography}{9}\small
\bibitem[Bratman(1992)]{Bratman1992} Bratman, M.~E. (1992). Practical reasoning and acceptance in a context. \emph{Mind}, 101(401), 1--15.
\bibitem[Fantl and McGrath(2002)]{FantlMcGrath2002} Fantl, J. and McGrath, M. (2002). Evidence, pragmatics, and justification. \emph{The Philosophical Review}, 111(1), 67--94.
\bibitem[Hawthorne(2004)]{Hawthorne2004} Hawthorne, J. (2004). \emph{Knowledge and Lotteries}. Oxford University Press.
\bibitem[Hawthorne and Stanley(2008)]{HawthorneStanley2008} Hawthorne, J. and Stanley, J. (2008). Knowledge and action. \emph{The Journal of Philosophy}, 105(10), 571--590.
\end{thebibliography}

\end{document}
